\chapter{Combinatorial}

\section{Permutations}
	\subsection{Factorial}
		\import{factorial.tex}
		\kactlimport{IntPerm.h}

	\subsection{Cycles}
		Let $g_S(n)$ be the number of $n$-permutations whose cycle lengths all belong to the set $S$. Then
		$$\sum_{n=0} ^\infty g_S(n) \frac{x^n}{n!} = \exp\left(\sum_{n\in S} \frac{x^n} {n} \right)$$

	\subsection{Derangements}
		Permutations of a set such that none of the elements appear in their original position.
		\[ \mkern-2mu D(n) = (n-1)(D(n-1)+D(n-2)) = n D(n-1)+(-1)^n = \left\lfloor\frac{n!}{e}\right\rceil \]

	\subsection{Burnside's lemma}
		Given a group $G$ of symmetries and a set $X$, the number of elements of $X$ \emph{up to symmetry} equals
		 \[ {\frac {1}{|G|}}\sum _{{g\in G}}|X^{g}|, \]
		 where $X^{g}$ are the elements fixed by $g$ ($g.x = x$).

		 If $f(n)$ counts ``configurations'' (of some sort) of length $n$, we can ignore rotational symmetry using $G = \mathbb Z_n$ to get
		 \[ g(n) = \frac 1 n \sum_{k=0}^{n-1}{f(\text{gcd}(n, k))} = \frac 1 n \sum_{k|n}{f(k)\phi(n/k)}. \]

\section{Partitions and subsets}
	\subsection{Partition function}
		Number of ways of writing $n$ as a sum of positive integers, disregarding the order of the summands.
		\[ p(0) = 1,\ p(n) = \sum_{k \in \mathbb Z \setminus \{0\}}{(-1)^{k+1} p(n - k(3k-1) / 2)} \]
		\[ p(n) \sim 0.145 / n \cdot \exp(2.56 \sqrt{n}) \]

		\begin{center}
		\begin{tabular}{c|c@{\ }c@{\ }c@{\ }c@{\ }c@{\ }c@{\ }c@{\ }c@{\ }c@{\ }c@{\ }c@{\ }c@{\ }c}
			$n$    & 0 & 1 & 2 & 3 & 4 & 5 & 6  & 7  & 8  & 9  & 20  & 50  & 100 \\ \hline
			$p(n)$ & 1 & 1 & 2 & 3 & 5 & 7 & 11 & 15 & 22 & 30 & 627 & $\mathtt{\sim}$2e5 & $\mathtt{\sim}$2e8 \\
		\end{tabular}
		\end{center}

	\subsection{Lucas' Theorem}
		Let $n,m$ be non-negative integers and $p$ a prime. Write $n=n_kp^k+...+n_1p+n_0$ and $m=m_kp^k+...+m_1p+m_0$. Then $\binom{n}{m} \equiv \prod_{i=0}^k\binom{n_i}{m_i} \pmod{p}$.

	\subsection{Binomials}
		\kactlimport{multinomial.h}

\section{General purpose numbers}
	\subsection{Bernoulli numbers}
		EGF of Bernoulli numbers is $B(t)=\frac{t}{e^t-1}$ (FFT-able).
		$B[0,\ldots] = [1, -\frac{1}{2}, \frac{1}{6}, 0, -\frac{1}{30}, 0, \frac{1}{42}, \ldots]$

		Sums of powers:
		\small
		\[ \sum_{i=1}^n i^m = \frac{1}{m+1} \sum_{k=0}^m \binom{m+1}{k} B_k \cdot (n+1)^{m+1-k} \]
		\normalsize

		Euler-Maclaurin formula for infinite sums:
		\small
		\[ \sum_{i=m}^{\infty} f(i) = \int_m^\infty f(x) dx - \sum_{k=1}^\infty \frac{B_k}{k!}f^{(k-1)}(m) \]
		\[ \approx \int_{m}^\infty f(x)dx + \frac{f(m)}{2} - \frac{f'(m)}{12} + \frac{f'''(m)}{720} + O(f^{(5)}(m)) \]
		\normalsize

	\subsection{Stirling numbers of the first kind}
		Number of permutations on $n$ items with $k$ cycles.
		\begin{align*}
			&c(n,k) = c(n-1,k-1) + (n-1) c(n-1,k),\ c(0,0) = 1 \\
			&\textstyle \sum_{k=0}^n c(n,k)x^k = x(x+1) \dots (x+n-1)
		\end{align*}
		$c(8,k) = 0, 5040, 13068, 13132, 6769, 1960, 322, 28, 1$ \\
		$c(n,2) = 0, 0, 1, 3, 11, 50, 274, 1764, 13068, 109584, \dots$

	\subsection{Eulerian numbers}
		Number of permutations $\pi \in S_n$ in which exactly $k$ elements are greater than the previous element. $k$ $j$:s s.t. $\pi(j)>\pi(j+1)$, $k+1$ $j$:s s.t. $\pi(j)\geq j$, $k$ $j$:s s.t. $\pi(j)>j$.
		$$E(n,k) = (n-k)E(n-1,k-1) + (k+1)E(n-1,k)$$
		$$E(n,0) = E(n,n-1) = 1$$
		$$E(n,k) = \sum_{j=0}^k(-1)^j\binom{n+1}{j}(k+1-j)^n$$

	\subsection{Stirling numbers of the second kind}
		Partitions of $n$ distinct elements into exactly $k$ groups.
		$$S(n,k) = S(n-1,k-1) + k S(n-1,k)$$
		$$S(n,1) = S(n,n) = 1$$
		$$S(n,k) = \frac{1}{k!}\sum_{j=0}^k (-1)^{k-j}\binom{k}{j}j^n$$

	\subsection{Bell numbers}
		Total number of partitions of $n$ distinct elements. $B(n) =$
		$1, 1, 2, 5, 15, 52, 203, 877, 4140, 21147, \dots$. For $p$ prime,
		\[ B(p^m+n)\equiv mB(n)+B(n+1) \pmod{p} \]

	\subsection{Labeled unrooted trees}
		\# on $n$ vertices: $n^{n-2}$ \\
		\# on $k$ existing trees of size $n_i$: $n_1n_2\cdots n_k n^{k-2}$ \\
		\# with degrees $d_i$: $(n-2)! / ((d_1-1)! \cdots (d_n-1)!)$

	\subsection{Pr\"ufer codes}
		Labeled trees on $n$ vertices $\leftrightarrow$ sequences in $\{1..n\}^{n-2}$ (a bijection, hence $n^{n-2}$ trees); vertex $v$ appears $\deg(v)-1$ times.
		Encode: repeatedly remove the smallest leaf and append its neighbour. Decode: the leaf removed at step $i$ is the smallest label that is not yet removed and does not occur in $a_i, \dots, a_{n-2}$; join it to $a_i$; finally join the two survivors. Both $O(n)$ with a pointer walk (\texttt{PruferCode.h}).
		Rooted forests with $k$ trees rooted at $1..k$: $k\,n^{n-k-1}$. Uniform random tree: random sequence, decode.
		\kactlimport{PruferCode.h}

	\subsection{Catalan numbers}
		\[ C_n=\frac{1}{n+1}\binom{2n}{n}= \binom{2n}{n}-\binom{2n}{n+1} = \frac{(2n)!}{(n+1)!n!} \]
		\[ C_0=1,\ C_{n+1} = \frac{2(2n+1)}{n+2}C_n,\ C_{n+1}=\sum C_iC_{n-i} \]
		${C_n = 1, 1, 2, 5, 14, 42, 132, 429, 1430, 4862, 16796, 58786, \dots}$
		\begin{itemize}
			\item sub-diagonal monotone paths in an $n\times n$ grid.
			\item strings with $n$ pairs of parenthesis, correctly nested.
			\item binary trees with with $n+1$ leaves (0 or 2 children).
			\item ordered trees with $n+1$ vertices.
			\item ways a convex polygon with $n+2$ sides can be cut into triangles by connecting vertices with straight lines.
			\item permutations of $[n]$ with no 3-term increasing subseq.
		\end{itemize}
\section{Generating functions}
	OGF $F=\sum f_nx^n$, EGF $\hat F=\sum f_nx^n/n!$; formal, convergence never matters. \textbf{Which:} gluing sizes $k$, $n{-}k$ by a plain $\sum_kf_kg_{n-k}$ $\Rightarrow$ OGF (tiles, coins, paths, strings); by $\sum_k\binom nkf_kg_{n-k}$, i.e.\ the $n$ labelled atoms get split between the parts $\Rightarrow$ EGF (permutations, set partitions, labelled trees). For the equation, multiply the recurrence by $x^n$ and sum wherever it holds: $\sum_{n\ge k}a_{n-k}x^n=x^kA$, $\sum_nna_nx^n=xA'$, $\sum_n(\sum_{k<n}a_ka_{n-1-k})x^n=xA^2$. So a convolution makes it \emph{quadratic} (solve, keep the root with $A(0)$ finite: $a_n=\sum_{k<n}a_ka_{n-1-k}$, $a_0{=}1\Rightarrow A=1+xA^2$); an $na_n$ makes it an ODE.

	\subsection{Dictionary}
		\begin{tabular}{@{}l@{\ }l@{\ }l@{\ }l@{}}
			$F$ & $[x^n]F$ & $F$ & $[x^n]F$ \\ \hline
			$(1-x)^{-1}$ & $1$ & $(1-ax)^{-1}$ & $a^n$ \\
			$(1-ax)^{-k}$ & $\binom{n+k-1}{k-1}a^n$ & $(1+x)^r$ & $\binom rn$ \\
			$x(1-x)^{-2}$ & $n$ & $x^k(1-x)^{-k-1}$ & $\binom nk$ \\
			$\ln\frac1{1-x}$ & $1/n$, $n\ge1$ & $\frac{x(1+x)}{(1-x)^3}$ & $n^2$ \\
			$\frac{1-x^{m+1}}{1-x}$ & $[n\le m]$ & $\frac1{1-x}\ln\frac1{1-x}$ & $H_n$ \\
			$\frac x{1-x-x^2}$ & $F_n$ & $\frac{2-x}{1-x-x^2}$ & $L_n$ \\
			$(1-4x)^{-1/2}$ & $\binom{2n}n$ & $C^k/\sqrt{1-4x}$ & $\binom{2n+k}n$ \\
			$C(x)^k$, $k\ge1$ & $\frac k{2n+k}\binom{2n+k}n$ & $\frac{xA_k(x)}{(1-x)^{k+1}}$ & $n^k$, $k\ge1$ \\
			$\prod_k(1-x^k)^{-1}$ & $p(n)$ & $(1-x^m)^{-1}$ & $[m \mid n]$ \\ \hline
			\multicolumn{4}{@{}l@{}}{EGF: the entry is now $n![x^n]\hat F$.} \\
			$e^x$ & $1$ & $e^{ax}$ & $a^n$ \\
			$(1-x)^{-1}$ & $n!$ & $-\ln(1-x)$ & $(n-1)!$, $n\ge1$ \\
			$e^{-x}/(1-x)$ & $D_n$ & $e^{x+x^2/2}$ & involutions \\
			$(e^x-1)^k$ & $k!\,S(n,k)$ & $(-\ln(1-x))^k$ & $k!\,c(n,k)$ \\
			$e^{e^x-1}$ & Bell $B_n$ & $(2-e^x)^{-1}$ & ord.\ set part. \\
			$T=xe^T$ & $n^{n-1}$ & $\exp\sum_S\frac{x^s}{s!}$ & blocks in $S$ \\
		\end{tabular}

		$C=\frac{1-\sqrt{1-4x}}{2x}$ solves $C=1+xC^2$; $\sqrt{1-4x}=1-2\sum_{n\ge1}C_{n-1}x^n$; $F_0{=}0$, $L_0{=}2$; $A_k=\sum_{j<k}E(k,j)x^j$ is Eulerian ($A_3{=}1{+}4x{+}x^2$); $(e^x-1)^k$ also counts surjections onto $[k]$, $\exp\sum_Sx^s/s$ cycle lengths in $S$; as an OGF, $\sum_nS(n,k)x^n=\frac{x^k}{(1-x)\cdots(1-kx)}$.

	\subsection{Operations}
		\begin{tabular}{@{}l@{\ }l@{\ \ }l@{\ }l@{}}
			$x^kF$ & $f_{n-k}$ & $F(ax)$ & $a^nf_n$ \\
			$F/(1-x)$ & $\sum_{k\le n}f_k$ & $(1-x)F$ & $f_n-f_{n-1}$ \\
			$xF'$ & $nf_n$ & $\int_0^xF$ & $f_{n-1}/n$ \\
			$FG$ & $\sum_kf_kg_{n-k}$ & $\hat F\hat G$ & $\sum_k\binom nkf_kg_{n-k}$ \\
		\end{tabular}

		Binomial transform $f_n\mapsto\sum_k\binom nkf_k$: OGF $\frac1{1-x}F(\frac x{1-x})$, EGF $e^x\hat F$, inverted by $e^{-x}\hat F$ (= ``forbid fixed points''). For $f_0{=}0$: \textbf{sequences} $\frac1{1-F}$ ($F^k$ = exactly $k$); labelled \textbf{sets} $e^{\hat F}$ ($\hat F^k/k!$), \textbf{cycles} $\ln\frac1{1-\hat F}$, so ``sets of connected pieces'' gives $C=\ln A$; unlabelled \textbf{multisets} $\exp\sum_{j\ge1}F(x^j)/j$, \textbf{distinct} $\exp\sum_{j\ge1}(-1)^{j-1}F(x^j)/j$.

	\subsection{Linear recurrences}
		$a_n=c_1a_{n-1}+\dots+c_ka_{n-k}$ ($n\ge k$) $\Rightarrow A=P/Q$, $Q=1-c_1x-\dots-c_kx^k$, $P=QA$ truncated below $x^k$ ($p_n=a_n-\sum_{j=1}^nc_ja_{n-j}$); conversely any $P/Q$ with $\deg P<\deg Q=k$ obeys it: rational $\Leftrightarrow$ linear recurrence. \textbf{Closed form}, if $c_k\ne0$ (else drop trailing zero $c_j$'s and shift): $\chi(t)=t^kQ(1/t)=t^k-c_1t^{k-1}-\dots-c_k$ is $Q$ reversed, and $\chi=\prod_i(t-r_i)^{m_i}$ gives $Q=\prod_i(1-r_ix)^{m_i}$ and $a_n=\sum_iq_i(n)r_i^n$, $\deg q_i<m_i$: \textbf{a root of multiplicity $m$ carries a degree-$(m{-}1)$ polynomial in $n$}. Fit the $k$ coefficients to $a_0..a_{k-1}$ by one linear solve, faster by hand than partial fractions ($[x^n](1-rx)^{-m}=\binom{n+m-1}{m-1}r^n$).

		\textbf{Inhomogeneous} $a_n=\sum_jc_ja_{n-j}+g(n)$, $g(n)=(\text{degree-}d$ poly$)\cdot b^n$ (polynomial: $b{=}1$; geometric: $d{=}0$; several $g$ superpose). Then $a$ obeys the \emph{homogeneous} recurrence with characteristic polynomial $\chi(t)(t-b)^{d+1}$, order $k{+}d{+}1$: compute that many terms, hand to \texttt{LinearRecurrence.h}. By hand, particular solution $n^mq(n)b^n$, $\deg q\le d$, $m$ = multiplicity of $b$ in $\chi$ ($0$ if $\chi(b)\ne0$); add the homogeneous part, fit. \emph{Ex.} $a_n=2a_{n-1}+n$, $a_0{=}0$: $\alpha n+\beta=-n-2$, $a_n=2^{n+1}-n-2$. Prefix sums are $\chi(t)(t-1)$, i.e.\ $A/(1-x)$.

		$a_n=\alpha_na_{n-1}+\beta_n\Rightarrow a_n=A_n(a_0+\sum_{j=1}^n\beta_j/A_j)$, $A_n=\prod_{i=1}^n\alpha_i$, $A_0{=}1$. Non-constant polynomial $c_j$: no rational GF, but $nf_n\leftrightarrow xF'$ makes an ODE. $a_n$ polynomial in $n$ of degree $\le d$ $\iff A=P/(1-x)^{d+1}$, $\deg P\le d$ (\texttt{LagrangeInterpolation.h}). Any DP with a fixed state set and constant transitions $M$ is rational with denominator $\det(I-xM)$, so of order $\le$ \#states --- your licence to guess.

	\subsection{Extracting $[x^n]$}
		$(1+y)^r=\sum_{n\ge0}\binom rny^n$ for \emph{any} real $r$, $\binom rn=\frac{r(r-1)\cdots(r-n+1)}{n!}$, $=0$ for $n<0$; upper negation $\binom rn=(-1)^n\binom{n-r-1}n$ gives $[x^n](1-x)^{-k}=\binom{n+k-1}{k-1}$. Beware: ``$\binom nk=0$ for $k>n$'' needs $n$ a non-negative \emph{integer}. Sums become coefficients, $\sum_ka_kb_{n-k}=[x^n]AB$ (Vandermonde is $(1{+}x)^m(1{+}x)^n$).

		\textbf{Roots of unity filter}, $\omega$ of order $m$ ($e^{2\pi i/m}$; mod $p$ one exists iff $m\mid p-1$): $\sum_{n\equiv r\,(m)}f_nx^n=\frac1m\sum_{j<m}\omega^{-jr}F(\omega^jx)$, and $x=1$ sums those $f_n$. For $m=2$ the even part is the \emph{series} $\frac{F(x)+F(-x)}2$, the odd $\frac{F(x)-F(-x)}2$. So $\sum_{k\equiv r\,(m)}\binom nk=\frac1m\sum_j\omega^{-jr}(1+\omega^j)^n$; e.g.\ $\sum_k\binom n{2k}=2^{n-1}$ for $n\ge1$.

		\textbf{Lagrange inversion} for implicit $F=x\phi(F)$, $\phi(0)\ne0$: for $n\ge1$, $[x^n]F^k=\frac kn[u^{n-k}]\phi^n$ and $[x^n]H(F)=\frac1n[u^{n-1}]H'\phi^n$. \emph{Ex.} $xC$ obeys $F=x\frac1{1-F}$, so $C_{n-1}=\frac1n[u^{n-1}](1-u)^{-n}=\frac1n\binom{2n-2}{n-1}$; $T=xe^T$ gives $\frac1n[u^{n-1}]e^{nu}=n^{n-1}/n!$; $p$-ary trees $B=1+xB^p$ give $[x^n]B=\frac1{(p-1)n+1}\binom{pn}n$.

		\textbf{Rational $P/Q$}, $\deg Q=k$: every term up to $n$ in \bigo{nk} by running the recurrence, one far term in \bigo{k^2\log n} by \texttt{LinearRecurrence.h}.

	\subsection{If you see this}
		\textbf{$x_1+\dots+x_k=n$:} $x_i\ge0$: $\binom{n+k-1}{k-1}$; $x_i\ge1$: $\binom{n-1}{k-1}$; $0\le x_i\le c$: GF $(\frac{1-x^{c+1}}{1-x})^k$, count $\sum_{j\ge0}(-1)^j\binom kj\binom{n-j(c+1)+k-1}{k-1}$, dropping $j$ with $n<j(c{+}1)$; unequal $c_i$ or parity/step constraints: multiply the per-variable GFs.

		\textbf{Ordered parts} (compositions, $1\times n$ tilings, stair climbing) with sizes in $S$: $\frac1{1-\sum_{s\in S}x^s}$ --- the denominator \emph{is} the recurrence; exactly $k$ parts, $[x^n](\sum_sx^s)^k$. Any parts $\ge1$: $\frac{1-x}{1-2x}$, i.e.\ $2^{n-1}$ for $n\ge1$ ($1$ at $n{=}0$), exactly $k$ of them $\binom{n-1}{k-1}$. $S{=}\{1,2\}$: $F_{n+1}$, as are $2\times n$ domino tilings; wider strips: profile DP, then guess. Binary strings, no two adjacent 1s: $F_{n+2}$; no $k$ consecutive: $\frac{1-x^k}{1-2x+x^{k+1}}$. $k$ non-adjacent of $n$: $\binom{n-k+1}k$ in a row, $\frac n{n-k}\binom{n-k}k$ in a circle ($0{<}k{<}n$); totals $F_{n+2}$, $L_n$.

		\textbf{Unordered parts} (partitions, change): parts from $S$, $\prod_{s\in S}(1-x^s)^{-1}$ (the knapsack DP). Distinct $\prod_k(1+x^k)$ $=$ odd $\prod_k(1-x^{2k-1})^{-1}$; parts $\le m$ $=$ at most $m$ parts $=\prod_{i\le m}(1-x^i)^{-1}$; exactly $k$ parts $x^k\prod_{i\le k}(1-x^i)^{-1}$, distinct and exactly $k$ $x^{\binom{k+1}2}\prod_{i\le k}(1-x^i)^{-1}$; each part $\le m$ times $\prod_k\frac{1-x^{k(m+1)}}{1-x^k}$; in a $k\times m$ box $\prod_{i=1}^k\frac{1-x^{m+i}}{1-x^i}$ (Gaussian binomial, degree $km$, sums to $\binom{k+m}k$).

		\textbf{Catalan-shaped:} anything with $A=1+xA^2$. Reflection: $\pm1$ steps, $n$ of them, ending at height $h$, never below $0$: $\binom nd-\binom n{d-1}$ with $d=\frac{n-h}2$ down-steps --- $0$ unless $d\in\mathbb Z$, $0\le d\le n/2$, and $\binom n{-1}{=}0$; no floor, $\binom nd$. Motzkin ($\pm1,0$, stay $\ge0$, end at $0$): $M=1+xM+x^2M^2$, $M_n=\sum_k\binom n{2k}C_k$. Narayana ($j$ peaks, semilength $n$): $\frac1n\binom nj\binom n{j-1}$, $\sum_j=C_n$. Steps E, N, NE: Delannoy $\frac1{1-x-y-xy}=\sum_j\binom mj\binom nj2^j$, or below the diagonal large Schr\"oder $\frac{1-x-\sqrt{1-6x+x^2}}{2x}$.

	\subsection{In a contest}
		Brute-force 30--40 exact terms mod the problem's prime and eyeball them: $1,2,5,14$ Catalan; $1,2,4,9,21$ Motzkin; $1,2,5,15,52$ Bell; $2,6,20,70$ $\binom{2n}n$; $1,3,13,75$ ordered set partitions; $1,2,4,10,26$ involutions; $1,2,9,44$ derangements; $3,13,63,321$ Delannoy; $2,6,22,90$ Schr\"oder. Then give $\ge2k$ of them to \texttt{BerlekampMassey.h}; if a short \texttt{tr} comes back and predicts terms you withheld, \texttt{LinearRecurrence.h} finishes it. Both use $\texttt{tr}=(c_1..c_k)$, i.e.\ $Q=1-\sum_jc_jx^j$. \emph{Guessing beats deriving} and works whenever the GF is rational --- every tiling, path, automaton or fixed-state-DP count. Only then do the algebra: a $\sqrt{\ }$ or $\exp$ in $A$ means no constant-coefficient recurrence exists (try Lagrange inversion); terms like $n!$ mean retry with the EGF.
